\section{All-prime Fermat sections at regular exponents}
\label{sec:allprime-regular-bridge}

One global affine unit turns the regular-primary local cusp condition
into a common centre at every finite place. Write
\[
 F_n:\ X^n+Y^n=Z^n,\qquad
 U_n=F_n\setminus\{XYZ=0\},\qquad
 \mathcal D_n=F_n\setminus U_n.
\]
Throughout this section \(n\geq3\) is odd.  The three rational
coordinate points are
\[
 P_0=(0:1:1),\qquad P_1=(1:0:1),\qquad
 P_\infty=(1:-1:0).
\]
The notation \(\mathcal D_n\) includes all \(3n\) geometric boundary
points, not just these three rational points.  Changing the sign of
\(Z\) identifies this model with the all-plus Fermat model.

Put \(\Delta_n=\pi_1(U_{n,\overline{\mathbb Q}})\), and define
\[
 \Delta_n^{[0]}=\Delta_n,\qquad
 \Delta_n^{[d+1]}=
 \overline{[\Delta_n^{[d]},\Delta_n^{[d]}]},\qquad
 R_{n,d}=\Delta_n/\Delta_n^{[d]}\quad(d\geq1).
\]
Also put
\[
 R_{n,\mathrm{sol}}
 =\Delta_n\big/\bigcap_{d\geq0}\Delta_n^{[d]}.
\]
Every arithmetic quotient considered below retains the entire group
\(G_{\mathbb Q}\).  Section classes are taken modulo conjugation by
the indicated geometric group.  The local condition means that,
at each finite place, the quotient section is the image of either
a full local point section of \(U_n\), or a full section of the
decomposition group of a \(\mathbb Q_v\)-rational point of
\(\mathcal D_n\).  At the real place the analogous point-or-cusp
condition is imposed.  A full global lift is not part of this
definition.

\begin{theorem}[All-prime cusp packets at regular exponents]
\label{thm:allprime-regular-bridge}
Suppose that \(n\) is divisible by \(3\), by \(5\), or by an odd
regular prime \(p>5\).  For every \(d\geq1\), and for the maximal
geometrically prosolvable quotient, every point-or-cusp Selmer
section is cuspidal.  After choosing one rational tangent at each
rational coordinate point, its section set is
\begin{equation}
 \coprod_{P\in\{P_0,P_1,P_\infty\}}
 H^1(\mathbb Q,\widehat{\mathbb Z}(1)).
 \label{eq:allprime-packets}
\end{equation}
Each displayed packet maps injectively to the section set, and
the three images are disjoint.  Every permitted full local
representative at a finite place has the same rational coordinate
centre as its global packet.

No common finite ramification set is required for the section or
its cusp parameter.
\end{theorem}

\subsection{Full Kummer completion and one fixed cusp}

We use the power completion
\begin{equation}
 \widehat{\mathbb Q^\times}
 =\varprojlim_{N\geq1}\mathbb Q^\times/(\mathbb Q^\times)^N
 =H^1(\mathbb Q,\widehat{\mathbb Z}(1)).
 \label{eq:allprime-power-completion}
\end{equation}
The inverse limit is over divisibility.  This notation does not
assert that the individual power quotients, or this inverse limit,
are compact.  The Kummer identification follows from the finite
Kummer sequences; the finite groups of rational roots of unity
satisfy the Mittag--Leffler condition on a cofinal factorial
subsystem.

\begin{lemma}[Valuations and local Kummer maps]
\label{lem:allprime-kummer}
The common kernel of all completed valuations
\[
 v_q:\widehat{\mathbb Q^\times}\longrightarrow\widehat{\mathbb Z}
 \qquad(q\text{ a rational prime})
\]
is \(\{1,-1\}\).  Multiplication by any odd positive integer is
injective on \(\widehat{\mathbb Q^\times}\).

For every finite place \(v\), the map from
\(\mathbb Q_v^\times\) to its full power completion is injective.
For an abelian variety \(A/\mathbb Q_v\), the map
\[
 A(\mathbb Q_v)\longrightarrow
 \varprojlim_N A(\mathbb Q_v)/N A(\mathbb Q_v)
\]
is injective as well.
\end{lemma}

\begin{proof}
Unique factorization gives
\[
 \mathbb Q^\times=\{1,-1\}\times
                     \bigoplus_q q^{\mathbb Z}.
\]
At every finite power level, a class whose valuation at each prime
is zero is represented by a sign.  Compatibility over the even
power levels makes that sign independent of the level.  This
proves the first assertion.

If \(\alpha^m=1\) and \(m\) is odd, then
\(m v_q(\alpha)=0\) for every \(q\).
The group \(\widehat{\mathbb Z}=\prod_\ell\mathbb Z_\ell\)
is torsion-free.  Thus all valuations vanish, so
\(\alpha\in\{1,-1\}\); oddness gives \(\alpha=1\).
In particular, for \(\epsilon\in\{1,-1\}\), the equation
\(\alpha^m=\epsilon\) has the unique solution
\(\alpha=\epsilon\).

The local multiplicative assertion follows from the usual
decomposition into valuation, finite roots of unity, and a
principal-unit group with an open \(\mathbb Z_v\)-lattice.
For the abelian assertion, \(A(\mathbb Q_v)\) is compact and has
an open subgroup \(\Lambda\simeq\mathbb Z_v^{\dim A}\);
see \citet[Lemma~I.3.3, p.~41]{Milne2006ADT}.
Choose a positive integer \(e\) annihilating
\(A(\mathbb Q_v)/\Lambda\).  If \(x\) is divisible by every
positive integer in \(A(\mathbb Q_v)\), then
\[
 x\in e v^r A(\mathbb Q_v)\subset v^r\Lambda
 \quad\text{for every }r,
\]
where \(v\) in this expression denotes the underlying rational
prime.  The intersection of these lattices is zero.
\end{proof}

In particular, full local Kummer cohomology distinguishes the
two rational signs.  A single odd-primary Kummer group need
not do so.  This distinction will be used at the prime \(p\).

\begin{lemma}[Recovering a full global cusp parameter]
\label{lem:allprime-one-cusp}
Let \(X/\mathbb Q\) be a smooth proper geometrically connected curve
of genus at least one, let \(D\) be a reduced finite divisor with at
least two geometric points, and put \(U=X\setminus D\).
Let \(P\in D(\mathbb Q)\).
Suppose a section of a finite full-geometric derived quotient, or
of the maximal geometrically prosolvable quotient, has at every
finite place a full local cusp representative at \(P\).
Then it is the image of a unique full global cusp section at \(P\),
where uniqueness refers to its
\(H^1(\mathbb Q,\widehat{\mathbb Z}(1))\) parameter after a
rational tangent has been chosen.
\end{lemma}

\begin{proof}
Choose a rational tangent at \(P\) and let \(t_P\) be its full section.
Filling \(P\) leaves an affine curve \(W=U\cup\{P\}\).
Lemma~\ref{lem:boundary} gives a unique
\(\beta\in H^1(\Q,\widehat\Z(1))\) whose peripheral image equals
the relative global abelian class of the given section.
Indeed, its image for \(W\) vanishes at every finite place,
and \citet[Proposition~3.2, author pp.~13--14]{Saidi2022} makes
that image globally zero. The lemma proves the global and local
peripheral injections, including their fixed-point hypothesis.

The parameter \(\beta\) defines a full cusp section \(t_{P,\beta}\).
At each finite place, the local peripheral injection identifies
its parameter with that of every full local cusp representative.
The two sections are therefore locally conjugate in the stated
quotient. Use the action from the constructed full section
\(t_{P,\beta}\), as required by the introduction of
\citet{Saidi2022}. His Proposition~3.1, author pp.~13--15,
makes the relative class globally zero at every derived depth
\(d\ge1\). His Theorem~B does the same in the prosolvable quotient.
Equivalently, the nonempty closed sets of conjugators giving
agreement at each derived depth are nested in the compact
geometric prosolvable group. A common conjugator gives agreement
in their inverse limit. This proves existence and uniqueness.
\end{proof}

The lemma imposes neither finite support on \(\beta\) nor a common
finite ramification set on its finite levels.  For example,
enumerate distinct rational primes \(q_j\), and specify completed
valuations \(v_{q_j}(\beta)=j!\), with all other valuations zero.
For each \(N\), only finitely many of these valuations are
nonzero modulo \(N\).  Their finite product defines a class in
\(\mathbb Q^\times/(\mathbb Q^\times)^N\), and these classes are
compatible as \(N\) varies.  They give a cusp parameter with
infinitely many nonzero valuations.

\subsection{The base regular prime}

\begin{proposition}[From one primary cusp to all finite centres]
\label{prop:allprime-base-regular}
Let \(p>5\) be regular.  Every section satisfying the local
condition on \(U_p\), in any quotient in
Theorem~\ref{thm:allprime-regular-bridge}, has all its full finite
local witnesses cuspidal at one rational coordinate point.
\end{proposition}

\begin{proof}
Choose arbitrary full local witnesses. Their projection to the
geometric pro-\(p\) abelian quotient satisfies the same local
condition. Corollary~\ref{cor:all-regular-primary-local} makes the
proper \(p\)-local centre a rational cusp \(P\).
Lemma~\ref{lem:all-regular-compact-boundary}, applied with \(n=p\),
then identifies the global compact \(T_pJ_p\)-class with that of
\(P\). At every \(q\ne p\), it also proves that a cusp witness
has centre \(P\), while an ordinary witness either reduces to
the open part or to \(P\). This use of the lemma concerns the
projection of the given quotient section and the same full
local witnesses; it assumes no full global lift.
A rational signed coordinate permutation lets us take
\(P=P_0=(0:1:1)\).
Put \(u=Y/Z\), a unit on \(U_p\) that extends as a unit at
\(P_0\), with \(u(P_0)=1\).

Pulling back the covers \(z^N=u\) along the quotient section
gives a single global class
\[
 \alpha\in H^1(\mathbb Q,\widehat{\mathbb Z}(1))
        =\widehat{\mathbb Q^\times}.
\]
All these covers have abelian geometric monodromy, so this
construction is available at every \(d\geq1\).
At a full local cusp section centred at \(P_0\), the value is
the Kummer class of \(u(P_0)=1\); inertia contributes nothing
because \(u\) extends as a unit there.
At an ordinary witness reducing to the open part, \(Y/Z\) is
a \(q\)-adic unit.
At an ordinary witness reducing to \(P_0\), both \(Y\) and \(Z\)
are units, so the same conclusion holds.
The \(p\)-local value is exactly \(1\).
We have proved
\[
 v_q(\alpha)=0\quad\text{for every rational prime }q,
 \qquad \alpha_p=1
 \quad\text{in }\widehat{\mathbb Q_p^\times}.
\]
Lemma~\ref{lem:allprime-kummer} first gives \(\alpha\in\{1,-1\}\).
Full local Kummer injectivity at \(p\) then gives \(\alpha=1\).

If a full witness at any finite \(q\) were an ordinary point
\(x_q\), the same local injectivity would imply
\(u(x_q)=1\).
The Fermat equation would then give \(X(x_q)=0\), contrary to
\(x_q\in U_p\).
Thus every finite witness is cuspidal at \(P_0\), as required.

The witnesses were arbitrary.  This proves the assertion for
every permitted choice, not just for a preferred system of
local representatives.
\end{proof}

Lemma~\ref{lem:allprime-one-cusp} completes the classification at
the base regular prime.

\subsection{Odd multiples without a good-reduction restriction}

\begin{proposition}[Odd-multiple transfer]
\label{prop:allprime-odd-transfer}
Let \(a\geq3\) be odd, and suppose that the conclusion of
Theorem~\ref{thm:allprime-regular-bridge} holds at exponent \(a\).
Then it holds at every exponent \(n=am\) with \(m\) odd.
\end{proposition}

\begin{proof}
The coordinate-power morphism
\[
 \phi_m:F_n\longrightarrow F_a,\qquad
 (X:Y:Z)\longmapsto(X^m:Y^m:Z^m)
\]
takes \(U_n\) to \(U_a\) and its boundary to the boundary.
Functoriality induces maps of all the stated finite-derived
arithmetic quotients and of their prosolvable limits.
The full arithmetic quotient \(G_{\mathbb Q}\) is unchanged.
There is no need to assert that a finite-derived quotient of a
covering subgroup embeds in the corresponding quotient below.

Push forward the section and each of its full local witnesses.
The resulting section at exponent \(a\) is globally cuspidal at
one rational coordinate point \(P_a\).
At each finite place \(v\), the proper full Tate-module Kummer
class of the image witness equals that of \(P_a\).
The Abel--Jacobi map is injective on \(F_a(\mathbb Q_v)\):
two distinct points with the same divisor class would give a
degree-one function on a curve of positive genus.
Its composition with full local Kummer completion remains
injective by Lemma~\ref{lem:allprime-kummer}.
Thus the image witness has proper centre exactly \(P_a\).
An ordinary point of \(U_n\) maps to an ordinary point of
\(U_a\); therefore every original finite witness is a cusp
witness at some
\[
 Q_v\in\mathcal D_n(\mathbb Q_v),\qquad \phi_m(Q_v)=P_a.
\]
This argument also applies at primes dividing \(m\).
It does not use good reduction there.

Choose signed rational coordinate permutations to make
\(P_a=P_0\).
Because \(m\) is odd, these permutations commute with the
coordinate-power maps.
Let \(u_n=Y/Z\) and \(u_a=Y/Z\).
Then \(\phi_m^*u_a=u_n^m\), and the \(u_a\)-Kummer class of the
pushed section is \(1\), since \(u_a\) extends as a unit of
value \(1\) at \(P_0\).
Consequently the global \(u_n\)-class satisfies
\[
 \alpha_n^m=1\quad\text{in }\widehat{\mathbb Q^\times}.
\]
Odd multiplication is injective, so \(\alpha_n=1\).

At every \(Q_v\) over \(P_0\), the function \(u_n\) extends as
a unit.  Full local Kummer injectivity gives \(u_n(Q_v)=1\).
Since \(X(Q_v)=0\), this says
\(Q_v=(0:1:1)\), the rational cusp \(P_0\) of \(F_n\).
Lemma~\ref{lem:allprime-one-cusp} now identifies the original
section with its unique full global cusp lift.
\end{proof}

\begin{proof}[Proof of Theorem~\ref{thm:allprime-regular-bridge}]
At a regular prime \(p>5\), use
Proposition~\ref{prop:allprime-base-regular} and
Lemma~\ref{lem:allprime-one-cusp}.
For the bases \(3\) and \(5\), the established arithmetic
inputs give
\begin{equation}
 F_a(\mathbb A_{\mathbb Q})_\bullet^{f\text{-}\mathrm{ab}}
       =F_a(\mathbb Q)
       =\{P_0,P_1,P_\infty\}\qquad(a=3,5).
 \label{eq:allprime-small-base-input}
\end{equation}
For \(a=3\), use the elliptic modular model \(X_0(27)\) and
\citet[Corollary~8.8, together with Theorem~8.6]{Stoll2007}.
For \(a=5\), the nonconstant map to the rank-zero abelian
variety with finite Tate--Shafarevich group in
Proposition~\ref{prop:arithmetic} supplies
\citet[Theorem~8.6, p.~385]{Stoll2007}.
The rational point lists and the modified real-place
convention are those established with these arithmetic
inputs; Stoll's real-place erratum is included.

For completeness, the step from
\eqref{eq:allprime-small-base-input} to sections uses no global
full lift.
A quotient section retaining proper geometric abelianization
makes its proper local centres survive every finite abelian
torsor: pull the torsor back to the section and twist by that
global class; the full local witnesses provide local lifts.
The displayed equality therefore makes all finite centres one
rational cusp.
Lemma~\ref{lem:allprime-one-cusp} gives the required full cusp
parameter and the section equality.
Proposition~\ref{prop:allprime-odd-transfer} now treats every
odd multiple in the theorem.

Every element of
\(H^1(\mathbb Q,\widehat{\mathbb Z}(1))\) defines a full
global cusp section, hence satisfies all the required local
conditions, including the real one.
Within one centre, injectivity follows from
Lemma~\ref{lem:boundary}.
If two different rational centres gave conjugate quotient
sections, their proper full Tate-module Kummer classes would
coincide at any finite place.
Local Kummer injectivity and injectivity of the Abel--Jacobi
map would make the two points equal, a contradiction.
This proves injectivity and disjointness of all three packets,
and hence Theorem~\ref{thm:regular-solvable-intro}.

No density or integrality condition was added to the local
hypothesis.  The only density-one theorem used is Sa\"idi's
localization theorem, applied after equality at every finite
place has already been proved.
\end{proof}

\begin{remark}[Scope of the bridge]
\label{rem:allprime-bridge-scope}
The proof uses all primes in geometric abelianization twice:
to select the rational sign of a global affine unit, and to
recover the full inertia parameter at a fixed cusp.
It does not assert diagonality of the finite-abelian-descent adelic set of the
proper curve at every regular exponent.
It is a theorem for the specified finite-derived and
prosolvable arithmetic quotients, with full local witnesses;
it does not deduce such a theorem for an arbitrary geometric
quotient or for a single-primary affine quotient.
\end{remark}
